\section{Introduction}\label{sec:intro}

Kubernetes clusters run many containers that share one kernel, and their runtime security in practice rests on rules:
Falco~\cite{falco}, Tetragon~\cite{tetragon} and Tracee~\cite{tracee} watch system calls and kernel functions through
eBPF and raise an alert when an event matches a pattern someone wrote. Rules are precise and cheap, but they only see
what their authors anticipated. Anomaly detection from system calls promises the complement, flagging behaviour a
workload never showed before, and the idea is thirty years old~\cite{forrest1996}. Yet most of that line of work is
evaluated offline, on recorded traces, often with thresholds tuned on labelled attacks~\cite{arp2022}. Three
questions that decide whether such a detector can run next to the rules on a live cluster are rarely measured: how long
an alert takes from the kernel event to the person who must act, how many alarms a real workload produces per hour of
normal operation, and whether the detector can tell an attack from a new but harmless process, such as an operator's
\texttt{kubectl exec}.

We present \sys, an eBPF-based detector that learns each workload's normal behaviour as a language over system-call
tokens. Each token keeps, besides the system-call name, the part of the return information that exploits change while
keeping the set of calls: the error code of a failure, the size class of a read or write, or the mode and first
directory of an opened file. The kernel builds one token per call, counts 4-token sequences per thread in a map, and
hands closed 0.5\,s windows to user space; a count language model with Witten--Bell smoothing scores each window by its
most surprising events, and a threshold fitted on held-out normal data only decides an alarm, which a bridge pushes
straight to Alertmanager.

We evaluate \sys under protocols registered before the data they concern was read, separating development results from
blind ones. Our contributions are:
\begin{itemize}
\item \textbf{In-kernel return-value tokens.} A probe that builds one 32-bit token per system call with its error code,
  returned size or opened directory, counts per-thread 4-grams in an eBPF map, and is checked by a proof of one token per
  call. It adds 0.4--0.5\,\textmu s per simple system call of a watched container (\cref{sec:eval-cost}).
\item \textbf{A label-free model that holds on datasets and live.} On the official split of LID-DS~2021 the final
  model detects 97.25\,\% of attack recordings with 0.19\,\% of normal recordings raising a false alarm and p95 latency
  1.54\,s; on CB-DS 891/891 breakouts. In four blind runs of 24 to 45 hours on a 21-workload application it detects every
  attack of the primary set; in the final 45-hour run 296/296 with no idle false alarm and p95 0.46\,s to the end of the
  alarm window, and every trial at 0.1 and 0.5 actions/s, on-off and padded with 300 seen calls per operation
  (\cref{sec:eval-datasets,sec:eval-live,sec:eval-final}).
\item \textbf{Which choices matter.} Returned sizes carry most of the signal; five of the 39 hooked system calls suffice
  on LID-DS; a moment Gumbel rule is not a correct distribution but the best normal-only threshold rule we tried; and
  the model beats a count rule on the same tokens at low false-alarm budgets (\cref{sec:eval-design}).
\item \textbf{Honest limits and a direction.} The model flags harmless new processes, is blind to attacks built from
  seen sequences (0/544), and ages after a cluster change (2.3$\times$ more threshold exceedances). We show that learning
  operator activity or gating the model does not fix the first limit, and introduce a kernel lineage tag that separates a
  workload's own process tree from operator sessions (\cref{sec:lineage,sec:eval-benign}).
\end{itemize}
On LID-DS~2019 the earlier configuration detects 90.88\,\% of attack recordings at 0.44\,\% false alarms but misses
Heartbleed; the final configuration, chosen on a development part of that dataset, detects 97.15\,\% of the remaining
attack recordings at 0.46\,\% false alarms, and both miss the latency target because the first unseen behaviour comes
seconds after the dataset's timestamps. The final configuration doubles the live alert burden (0.36 episodes per
replica-hour), most of it in windows that hold the commands of Kubernetes exec probes.
